% This is samplepaper.tex, a sample chapter demonstrating the
% LLNCS macro package for Springer Computer Science proceedings;
% Version 2.25 of 2026/09/03
%
\documentclass[runningheads]{llncs}
%
\usepackage[T1]{fontenc}
% T1 fonts will be used to generate the final print and online PDFs,
% so please use T1 fonts in your manuscript whenever possible.
% Other font encondings may result in incorrect characters.
%

\usepackage{newtxtext}       %
\usepackage[varvw]{newtxmath}       % selects Times Roman as basic font


\usepackage{graphicx}
% Used for displaying a sample figure. If possible, figure files should
% be included in EPS format.
%
% If you use the hyperref package, please uncomment the following two lines
% to display URLs in blue roman font according to Springer's eBook style:
%\usepackage{color}
%\renewcommand\UrlFont{\color{blue}\rmfamily}
%\urlstyle{rm}
%
\begin{document}
%
\title{BidBots: Auction-Based Coordination of Warehouse Robots}
\titlerunning{BidBots}

\author{Acalfoaie Bogdan-Mihai}
%\authorrunning{B.-M. Acalfoaie}

\institute{Faculty of Computer Science, Alexandru Ioan Cuza University, Iași, Romania\\
\email{bogdan-mihai.acalfoaie@info.uaic.ro}}

\maketitle

\begin{abstract}
This report presents BidBots, a client/server system that coordinates a
fleet of simulated warehouse robots. Tasks are allocated through auctions
based on the Contract Net Protocol, while collisions are prevented by having
robots reserve their paths in space and time. We describe the system's
architecture, its TCP-based application protocol, and usage scenarios
that cover both normal operation and robot failures.

\keywords{Computer networks \and TCP \and Contract Net Protocol \and Multi-robot path planning}
\end{abstract}

\section{Introduction}
Modern warehouses rely on fleets of autonomous robots that move products from shelves to delivery stations. Coordinating such a fleet raises two problems: deciding which robot should handle each order, balancing efficiency against battery levels, and making sure that robots sharing the same aisles never collide. BidBots is a client/server system that addresses both problems by coordinating a fleet of simulated robots on a grid-based warehouse map.

Tasks are assigned through auctions based on the Contract Net Protocol~\cite{cnp}: the server announces each delivery, the available robots estimate their own cost, and the task goes to the best bid. The main objectives of the project are:

\begin{itemize}
  \item \textbf{Auction-based task allocation}: each robot bids according to its distance and battery level.
  \item \textbf{Collision-free movement}: robots reserve the cells of their path in space and time before moving.
  \item \textbf{Fault tolerance}: robots that stop responding are detected through heartbeats, and their tasks and reservations are reassigned.
  \item \textbf{Battery awareness}: robots only bid on tasks they can complete and return to a charging station when their battery runs low.
  \item \textbf{Evaluation}: a comparison of the auction mechanism against a simple nearest-free-robot assignment.
\end{itemize}

\section{Applied Technologies}

\begin{itemize}
  \item \textbf{Transport layer protocol}: the server communicates with all clients over TCP, which guarantees reliable, in-order delivery. Losing or reordering a message (e.g., a task award or a reservation reply) could lead to two robots holding the same task or entering the same cell.

  \item \textbf{Concurrency with one thread per connection}: the main thread accepts incoming connections and creates a dedicated thread for each client. These client threads read their client's messages and update the shared server state accordingly. A separate coordinator thread handles time-based events, such as closing auctions and detecting robots that have stopped responding. All threads access the shared state through a mutex.
\end{itemize}

\section{Application Structure}

\subsection{Modeling Concepts}

\begin{itemize}
  \item \textbf{Warehouse map}: a 2D grid of cells, each being a corridor, a
        shelf, a charging station or a delivery station. In each step, a robot
        moves to a neighboring cell or waits in place.
  \item \textbf{Simulation clock}: time is divided into discrete steps
        (ticks), advanced by the organizer thread. All movements and
        reservations refer to these steps.
  \item \textbf{Robot}: identified by an ID, with a position, a battery level
        and a state (idle, executing a task, charging or dead).
  \item \textbf{Task}: a delivery from a shelf to a delivery station. A task
        goes through the states \emph{announced}, \emph{awarded},
        \emph{picked up} and \emph{delivered}, and is auctioned again if its
        robot fails.
  \item \textbf{Auction}: opened by the organizer for each task; it collects
        bids until a deadline, after which the task is awarded to the lowest
        bid.
  \item \textbf{Reservation}: a robot's claim on cell $(x, y)$ at step $t$.
        Robots plan their paths locally and must reserve them before moving;
        the server rejects conflicting reservations, including two robots
        swapping cells in the same step.
  \item \textbf{Heartbeat}: a periodic message through which a robot signals
        that it is alive. A robot that misses heartbeats for longer than a
        threshold is considered dead.
\end{itemize}

\subsection{Components}

Fig.~\ref{fig:arch} shows the architecture of the system. It consists of a
concurrent server and three types of clients, each running as a separate
process and communicating with the server over TCP:

\begin{itemize}
  \item \textbf{Robot client}: simulates a single robot. It bids on tasks,
        plans its path, reserves it, moves and sends heartbeats.
  \item \textbf{Order client}: submits delivery orders and receives
        confirmations.
  \item \textbf{Visual client}: subscribes to the server and displays the
        warehouse map and the robots in real time.
\end{itemize}

Inside the server, the main thread accepts incoming connections and spawns a
dedicated handler thread for each client. The organiser thread advances the
simulation clock, manages auctions and detects failed robots. All threads
communicate through the shared state, which is protected by a mutex.

\begin{figure}
\centering
\includegraphics[width=\textwidth]{architecture.png}
\caption{BidBots architecture. All clients first connect through the main
thread, which spawns a dedicated handler thread for each connection.}
\label{fig:arch}
\end{figure}

\section{Communication Protocol}
TODO

\section{Usage Scenarios}
\subsection{Success Scenarios}
TODO
\subsection{Failure Scenarios}
TODO

\section{Conclusions}
TODO

\begin{thebibliography}{8}
\bibitem{cnp}
Smith, R.G.: The Contract Net Protocol: High-Level Communication and Control
in a Distributed Problem Solver. IEEE Transactions on Computers C-29(12), 1104--1113 (1980)
\end{thebibliography}

\end{document}